\documentclass[10pt,a4paper]{amsart}
\usepackage[margin=19mm]{geometry}
\usepackage{amsmath,amssymb,mathtools}
\usepackage[scr=rsfs,cal=euler]{mathalfa}
\usepackage{xfrac}
\usepackage[unicode,hidelinks,bookmarks=false]{hyperref}
\newcommand{\ie}{i.e.\ }
\newcommand{\cf}{cf.\ }
\newcommand{\resp}{respectively\ }
\newcommand{\Subsection}{\subsection}
\providecommand{\nobreakdash}{\nobreak\hbox{-}}
\setlength{\emergencystretch}{2em}
\multlinegap0pt
\usepackage{amssymb}
\usepackage{enumerate}
\usepackage{tikz}
\usepackage{float}
\newtheorem{proposition}{Proposition}
\newtheorem{proofx}{Proofx}
\newcommand{\R}{\mathbb{R}}
\newcommand{\Z}{\mathbb{Z}}
\newcommand{\dd}{\,{\rm d}}
\newcommand{\h}{{\rm h}}
\newcommand{\opdiv}{\operatorname{div}}
\newcommand{\oprot}{\operatorname{rot}}
\title{Existence of steady Navier-Stokes flows \\ exterior to an infinite cylinder}
\author{Mitsuo Higaki, Ryoma Horiuchi}
\date{Source: arXiv:2310.09752v3 (CC BY 4.0)}
\begin{document}
\maketitle

\noindent\emph{Source and licence.} Existence of steady Navier-Stokes flows \\ exterior to an infinite cylinder. Version arXiv:2310.09752v3 (CC BY 4.0), distributed by arXiv under the Creative Commons Attribution 4.0 International licence, http://creativecommons.org/licenses/by/4.0/, as recorded in the arXiv licence metadata for this exact version and shown on the abstract page https://arxiv.org/abs/2310.09752v3. Full licence notice: \texttt{licenses/arxiv-2310.09752-CC-BY-4.0.txt}.\par\smallskip

\noindent\textit{Editorial scope.} Counted unit: the article's proposition prop.LPh.zero (source0.tex lines 834--851). The article's complete original proof of it is delivered with it (source0.tex lines 867--1023, one \texttt{proof} environment of 157 lines, which the article places away from the statement). No mathematics is written, repaired or re-derived: the statement and the proof are the article's own lines, in the article's own notation, and no retained line is substituted or re-flowed.  Retained uncounted as the source-local context the proof needs: lines 540--549; lines 775--821; lines 811--814; lines 818--820.  Nothing was omitted from any retained span, so the package carries no omission record.  The retained lines cite 1 works, whose entries are transcribed verbatim from the article's own bibliography source in the e-print and delivered with short editorial labels recorded in the item's reference mapping.  No retained line contains a figure environment, an image-inclusion command or drawing code, so no image is delivered and the package declares no asset dependency.  Editorial formatting only, in addition: an AMS \texttt{amsart} preamble that restates the article's own declarations and macros used by the retained lines, namely the environments \texttt{proposition}, \texttt{proofx} on separate counters, together with the standard packages those lines need.  The retained environments print in this document as Proposition 1, Proofx 1 in order of appearance, whereas the article prints the same items with the numbering of its own full document.  The complete native source of the article as distributed in the arXiv e-print for this exact version is bundled unchanged as \texttt{sources/148484/source0.tex}.
\begin{equation}\tag{LP$_\h$}\label{eq.LPh}
\left\{
\begin{array}{ll}
-\Delta_\h v_\h + (V_\h)^\bot \oprot_{{\rm 2D}} v_\h + \nabla_\h q 
= f_\h&\mbox{in}\ \Omega \\
\opdiv_\h v_\h = 0&\mbox{in}\ \Omega \\
v_\h = 0 &\mbox{on}\ \partial\Omega \\
v_\h(x)\to0&\mbox{as}\ |x|\to\infty 
\end{array}\right.
\end{equation}

Let $n\in\Z$. Applying ${\mathcal P}_n$ to $v_\h$ and $f_\h$
%
\begin{align*}
({\mathcal P}_n v_\h)(r,\theta)
&= v_{r,n}(r) e^{in\theta} {\bf e}_r 
+ v_{\theta,n}(r) e^{in\theta} {\bf e}_\theta, \\
({\mathcal P}_n f_\h)(r,\theta)
&= f_{r,n}(r) e^{in\theta} {\bf e}_r 
+ f_{\theta,n}(r) e^{in\theta} {\bf e}_\theta 
\end{align*}
%
and inserting these to \eqref{eq.LPh}, we see that $(v_{r,n}(r),v_{\theta,n}(r))$ and $q_n(r)$ satisfy 
%
\begin{align}
&\begin{aligned}\label{eq.LPh.polar.vr}
&-\frac{\dd}{\dd r} \Big(\frac1r\frac{\dd}{\dd r} (r v_{r,n})\Big)
+ \frac{n^2}{r^2} v_{r,n}
+ \frac{2in}{r^2} v_{\theta,n} \\
&\qquad\qquad
- \frac{\alpha}{r^2} \Big(\frac{\dd}{\dd r} (r v_{\theta,n}) - in v_{r,n}\Big)
+ \partial_r q_n 
= f_{r,n}, \quad r>1, \\
\end{aligned} \\
&\begin{aligned}\label{eq.LPh.polar.vtheta}
&-\frac{\dd}{\dd r} \Big(\frac1r\frac{\dd}{\dd r} (r v_{\theta,n})\Big) 
+ \frac{n^2}{r^2} v_{\theta,n}
- \frac{2in}{r^2} v_{r,n} \\
&\qquad\qquad
- \frac{\gamma}{r^2} \Big(\frac{\dd}{\dd r} (r v_{\theta,n}) - in v_{r,n}\Big)
+ \frac{i n}{r} q_n 
= f_{\theta,n}, \quad r>1, 
\end{aligned}
\end{align}
%
the divergence-free and boundary conditions 
%
\begin{align}\label{eq.LPh.polar.div-free.bdry}
\frac{\dd}{\dd r}(r v_{r,n}) + in v_{\theta,n}=0, 
\qquad v_{r,n}(1)=v_{\theta,n}(1)=0, 
\end{align}
%
and the condition at infinity
%
\begin{align}\label{eq.LPh.polar.spatial.infinity}
|v_{r,n}(r)| + |v_{\theta,n}(r)| \to 0, \quad r\to\infty. 
\end{align}
%

\begin{align}
\frac{\dd}{\dd r}(r v_{r,n}) + in v_{\theta,n}=0, 
\qquad v_{r,n}(1)=v_{\theta,n}(1)=0, 
\end{align}

\begin{align}
|v_{r,n}(r)| + |v_{\theta,n}(r)| \to 0, \quad r\to\infty. 
\end{align}

\begin{proposition}\label{prop.LPh.zero}
Let $\alpha\in\R$, $\gamma>2$ and $2<\rho\le\gamma$. Suppose that $f_\h=f_{\h,0}$ is given by $f_{\h,0}=\opdiv_\h F_0$ for some $F_0\in {\mathcal P}_0 L^\infty_{2(\rho-1)}(\Omega)^{2\times2}$. Then there is a unique weak solution $v_{\h,0}\in {\mathcal P}_0 L^2_\sigma(\Omega) \cap W^{1,2}_0(\Omega)^2$ of \eqref{eq.LPh} satisfying 
\[
v_{\h,0}(r,\theta)=v_{\theta,0}(r) {\bf e}_\theta
\]
and 
%
\begin{align}\label{est1.prop.LPh.zero}
\|v_{\h,0}\|_{L^\infty_{\rho-1}} 
+ \frac{1}{\gamma-1} 
\|\nabla_\h v_{\h,0}\|_{L^\infty_{\rho}} 
&\le 
\frac{C(\gamma-1)}{(\gamma-2)(\rho-2)} 
\|F_0\|_{L^\infty_{2(\rho-1)}}. 
\end{align}
%
The constant $C$ is independent of $\alpha$, $\gamma$ and $\rho$.
\end{proposition}

\begin{proofx}{Proposition \ref{prop.LPh.zero}}
To simplify, we omit the subscript ``$\h$" for $v_\h$ and $f_\h$. Put $n=0$ in \eqref{eq.LPh.polar.vr}--\eqref{eq.LPh.polar.spatial.infinity}. The radial part $v_{r,0}(r)$ is identically zero since both $(r v_{r,0})'=0$ and $v_{r,0}(1)=0$ hold by \eqref{eq.LPh.polar.div-free.bdry}. Besides, the angular part $v_{\theta,0}(r)$ solves the ordinary differential equation 
%
\begin{align}\label{eq1.proof.prop.LPh.zero}
-\frac{\dd^2 v_{\theta,0}}{\dd r^2} 
-\frac{1+\gamma}r \frac{\dd v_{\theta,0}}{\dd r}  
+\frac{1-\gamma}{r^2} v_{\theta,0}
= f_{\theta,0}, \quad r>1, 
\qquad 
v_{\theta,0}(1)=0. 
\end{align}
%
As the uniqueness of solutions is clear, we may only consider the existence and estimates. 


At first we assume that $F_0\in {\mathcal P}_0 C^\infty_0(\Omega)^{2\times2}$, which gives 
%
\begin{align}\label{eq2.proof.prop.LPh.zero}
f_{\theta,0}
= (\opdiv_\h F)_{\theta,0}
= \frac{1}{r} \frac{\dd}{\dd r} (r F_{r \theta,0}) 
+ \frac{1}{r} F_{\theta r,0}. 
\end{align}
%
As is shown in \cite[Proof of Proposition 3.1]{Higaki2023}, the solution of \eqref{eq1.proof.prop.LPh.zero} is represented by 
%
\begin{equation*}
\begin{split}
v_{\theta,0}(r)
&= 
\frac{1}{\gamma-2}
\bigg\{	
-\bigg(
\int_{1}^{\infty} s^2 f_{\theta,0}(s)\dd s 
\bigg) 
r^{-\gamma+1} \\
&\qquad\qquad\quad
+ r^{-\gamma+1} \int_{1}^{r} s^{\gamma} f_{\theta,0}(s)\dd s 
+ r^{-1} \int_{r}^{\infty} s^2 f_{\theta,0}(s)\dd s
\bigg\}. 
\end{split}
\end{equation*}
%
By integration by parts based on \eqref{eq2.proof.prop.LPh.zero}, we rewrite this formula as 
%
\begin{equation}\label{rep1.proof.prop.LPh.zero}
\begin{split}
&v_{\theta,0}(r) \\
&=\frac{1}{\gamma-2}
\bigg[
\bigg(
\int_{1}^{\infty} s F_{r \theta,0}(s) \dd s 
- \int_{1}^{\infty} s F_{\theta r,0}(s) \dd s 
\bigg) 
r^{-\gamma+1} \\
&\qquad\qquad\quad
+ r^{-\gamma+1}
\bigg\{
-(\gamma-1) \int_{1}^{r} s^{\gamma-1} F_{r \theta,0}(s) \dd s 
+ \int_{1}^{r} s^{\gamma-1} F_{\theta r,0}(s) \dd s 
\bigg\} \\
&\qquad\qquad\quad
+ r^{-1} 
\bigg(
-\int_{r}^{\infty} s F_{r \theta,0}(s) \dd s 
+ \int_{r}^{\infty} s F_{\theta r,0}(s) \dd s 
\bigg)
\bigg]. 
\end{split}
\end{equation}
%
An associated pressure ${\mathcal P}_0 q$ is obtained from the equation \eqref{eq.LPh.polar.vr}. Set $v_0(r,\theta)=v_{\theta,0}(r) {\bf e}_\theta$. Then the pair $(v_0, \nabla_\h{\mathcal P}_0 q)$ is smooth and satisfies \eqref{eq.LPh.polar.vr}--\eqref{eq.LPh.polar.spatial.infinity} in the classical sense. 


Next we let $F_0\in {\mathcal P}_0 L^\infty_{2(\rho-1)}(\Omega)^{2\times2}$. By direct computation 
%
\begin{align*}
\begin{split}
\bigg(\int_{1}^{\infty} 
s |F_0(s)| \dd s 
\bigg) 
r^{-\gamma+1} 
\le 
\frac{C}{\rho-2} 
\|F_0\|_{L^\infty_{2(\rho-1)}} 
r^{-\gamma+1} 
\end{split}
\end{align*}
%
and 
%
\begin{align*}
\begin{split}
&r^{-\gamma+1} 
\int_1^r 
s^{\gamma-1} |F_0(s)| \dd s 
+ r^{-1} 
\int_{r}^{\infty}
s |F_0(s)| \dd s 
\le 
\frac{C}{\rho-2} 
\|F_0\|_{L^\infty_{2(\rho-1)}} 
r^{-\rho+1}, 
\end{split}
\end{align*}
%
one can verify that the desired estimate 
\eqref{est1.prop.LPh.zero} follows from the formula \eqref{rep1.proof.prop.LPh.zero}. 


Finally, we show that $v_0(r,\theta)=v_{\theta,0}(r) {\bf e}_\theta$ with $v_{\theta,0}(r)$ defined in \eqref{rep1.proof.prop.LPh.zero} is a weak solution of \eqref{eq.LPh} for general $F_0\in {\mathcal P}_0 L^\infty_{2(\rho-1)}(\Omega)^{2\times2}$. By the definition of $v_0$, we have $v_0\in L^2_{\sigma}(\Omega)\cap W^{1,2}_0(\Omega)^2$. Let $\{F^{(m)}_0\}_{m=1}^\infty\subset {\mathcal P}_0 C^\infty_0(\Omega)^{2\times2}$ be such that $\displaystyle{\lim_{m\to\infty} F^{(m)}_0 = F_0}$ in $L^1(\Omega)\cap L^2(\Omega)$. Let $v^{(m)}_0$ be given by \eqref{rep1.proof.prop.LPh.zero} replacing $F_0$ by $F^{(m)}_0$, and let ${\mathcal P}_0 q^{(m)}$ be the associated pressure. Then, using the estimate for $G\in L^1(0,\infty; s\dd s)^{2\times2}$ 
%
\begin{align*}
r^{-\gamma} 
\int_1^r 
s^{\gamma-1} |G(s)| \dd s 
+ r^{-2} 
\int_{r}^{\infty}
s |G(s)| \dd s 
\le
\bigg(\int_{1}^{\infty} |G(s)| s\dd s\bigg) r^{-2}
\end{align*}
%
which implies 
\[
\|\nabla_\h (v_0-v^{(m)}_0)\|_{L^2} 
\le 
C(\|F_0-F^{(m)}_0\|_{L^1} + \|F_0-F^{(m)}_0\|_{L^2}),
\]
we have, by integration by parts, 
%
\begin{align*}
\begin{split}
&\langle \nabla_\h v_0, \nabla_\h \varphi \rangle 
+ \langle V_\h^\bot \oprot_{{\rm 2D}} v_0, \varphi \rangle 
+ \langle F_0, \nabla_\h \varphi \rangle \\
&= 
\lim_{m\to\infty}
\Big(
\langle \nabla_\h v^{(m)}_0, \nabla_\h \varphi \rangle
+ \langle V_\h^\bot \oprot_{{\rm 2D}} v^{(m)}_0, \varphi \rangle
+ \langle F^{(m)}_0, \nabla_\h \varphi \rangle
\Big) \\
&= \lim_{m\to\infty} 
\langle 
-\Delta_\h v^{(m)}_0 
+ V_\h^\bot \oprot_{{\rm 2D}} v^{(m)}_0 
- \opdiv_\h F^{(m)}_0, \varphi \rangle \\
&= \lim_{m\to\infty} 
\langle \nabla_\h {\mathcal P}_0 q^{(m)}, \varphi \rangle = 0, 
\quad 
\varphi\in C^\infty_{0,\sigma}(\Omega). 
\end{split}
\end{align*}
%
Hence we see that $v_0$ is a weak solution of \eqref{eq.LPh}. This completes the proof. 
\end{proofx}

\begin{thebibliography}{99}
\bibitem[Higaki]{Higaki2023} Mitsuo Higaki. \newblock Existence of planar non-symmetric stationary flows with large flux in an exterior disk. \newblock {\em J. Differential Equations}, 360:182--200, 2023.
\end{thebibliography}
\end{document}
